\section{A complete positive LC prefix for the quadratic tribonacci inverse}
Let $\lambda\in(1,2)$ be the positive root of
$\lambda^3-\lambda^2-\lambda-1=0$, and put
\[
 a=\lambda-1,\quad b=\lambda^{-1}=a^2+a-1,\qquad
 1-z-z^2-z^3=(1-\lambda z)(1+az+bz^2).
\]
Thus $a^3+2a^2=2$, $83/100<a<84/100$, and $b>0$.
For $u,r\ge0$, $w=u+r\ge1$, set
\[
 E=(1+z)^u(1+z+z^2)^r,\quad D=u+2r,\qquad
 f_n=[z^n]\frac{E}{1+az+bz^2}.
\]
\begin{theorem}[All binary--ternary counts, one exact exception]
When $D\ge2$, the finite sequence $f_0,\ldots,f_{D-2}$ is
strictly positive and globally LC, except at $(u,r)=(0,2)$.
That exception has a strictly negative centre-one minor.
At $D=1$ the indicated prefix is empty.
\end{theorem}
The sequence is ZERO-EXTENDED after $D-2$. The last TWO ordinary
quotient coefficients are discarded; they can be negative, even
for the packet $[2][3]^2$. No positivity of the whole FPS is claimed.

\paragraph{Keep both terminal errors and their exact positive norm.}
Put $\epsilon_0=f_D$, $\epsilon_1=f_{D-1}$ and
\[
 Q_a(x,y)=x^2+axy+(a^2+a-1)y^2.
\]
For $E=1$ the terminal pair is $(1,0)$, using $f_{-1}=0$.
On multiplication by $[2]$ the pair changes to
\[
       ((1-a)x-by,\ x+y),\qquad Q_{\rm new}=a^2Q_a(x,y).
\]
On multiplication by $[3]$ it changes to
\[
 (2(1-a)x-b(1-a)y,\ (1-a)x+(1-b)y),
                  \qquad Q_{\rm new}=(1-a)Q_a(x,y).
\]
Both identities are exact real polynomial identities under
$a^3+2a^2=2$; the ordinary forward coefficient recurrence identifies
the displayed terminal pairs. Complete two-count induction gives
\[
             Q_a(\epsilon_0,\epsilon_1)
                              =(a^2)^u(1-a)^r.            \tag{1}
\]
This retains the phase of both nonreal roots; neither error is set
to zero or replaced by a one-dimensional sign assumption.

Since $a^2<3/4$ and $0<1-a<3/4$, (1) is at most $(3/4)^w$.
For all real x,y the bounds on a imply
\[
 Q_a(x,y)\ge(1-a/2)x^2+(a^2+a/2-1)y^2
                                      \ge(x^2+y^2)/10.
\]
Here the first inequality is $a(x+y)^2/2\ge0$; the second uses
$1-a/2\ge.58$ and $a^2+a/2-1\ge.1039$.
The exact rational inequality $(3/4)^{18}<1/160$ then proves
\[
        w\ge18\quad\Longrightarrow\quad
                       |\epsilon_0|,|\epsilon_1|\le1/4. \tag{2}
\]
Six Lean reports verify both full norm transformations, the entire
two-natural-count norm induction, coercivity and the uniform power/box
payment. The sequence-to-terminal identification is the written
coefficient calculation above.

\paragraph{Every new uncovered read is an explicit polynomial corner.}
Suppose the old prefix through D-2 is positive and LC and w>=18.
Write $x_k=f_{D-k}$. Palindromicity of E gives the COMPLETE backward
reads needed by either extension:
\[
\begin{aligned}
 x_0&=\epsilon_0,&x_1&=\epsilon_1,\\
 x_2&=\lambda(1-ax_1-x_0),&
 x_3&=\lambda(w-ax_2-x_1),\\
 x_4&=\lambda\{\tbinom w2+r-ax_3-x_2\},&
 x_5&=\lambda\{\tbinom w3+r(w-1)-ax_4-x_3\}.
\end{aligned}                                                   \tag{3}
\]
All indices exist since D>=w>=18. On binary extension the true new
prefix agrees with the positive LC convolution of the old truncated
prefix at every read through D-2. Its only additional kept coefficient
is $x_1+x_2$, and its only uncovered interior minor is
\[
                 (x_2+x_3)^2-(x_1+x_2)(x_3+x_4).             \tag{4}
\]
On ternary extension the two additional kept coefficients are
$x_1+x_2+x_3$ and $x_0+x_1+x_2$; the only uncovered minors are
\[
\begin{split}
 &(x_2+x_3+x_4)^2-(x_1+x_2+x_3)(x_3+x_4+x_5),\\
 &(x_1+x_2+x_3)^2-(x_0+x_1+x_2)(x_2+x_3+x_4).              \tag{5}
\end{split}
\]
Thus three value inequalities and THREE minor inequalities suffice;
no unspecified middle rank is passed to an absent old LC premise.

All six expressions are strictly positive throughout
$w\ge18$, $0\le r\le w$, and the entire terminal box (2).
Here is the complete exact certificate, with its unbounded meaning:
substitute
\[
 w=18+X,\quad r=\theta w,\quad
 \epsilon_i=-\tfrac14+\tfrac12\eta_i,
 \qquad X\ge0,\quad0\le\theta,\eta_0,\eta_1\le1.
\]
Expand in powers of X and, for each coefficient, in the tensor
Bernstein basis on the unit cube. There are respectively
30,54,30 coefficients for (4)--(5), and 4,3,4 for the three values:
125 coefficients in total. Every one is nonnegative in
$\mathbb Q(\lambda)$, and every coefficient at $X^0$ is strictly
positive. The complete triples are recorded in
tribonacci_inverse_complete_certificate_data.json.
Finite Bernstein basis functions are nonnegative and sum to one,
so these coefficient facts prove strict positivity at ALL real
points of the stated unbounded domain, including the box boundary.

The independent verifier checks the WHOLE polynomial identities:
each side has degree at most (4,2,2,2) in (X,theta,eta0,eta1),
and exact equality is checked at ALL points of the 5-by-3-by-3-by-3
tensor interpolation grid. These 135 values per identity determine
it uniquely, by repeated univariate interpolation. This is a complete
polynomial-identity check, not a bounded parameter experiment.
The scalar algebraic signs use a certified rational interval of
width $10^{-50}$ for lambda; a complete Lean integer certificate
proves all required lower bounds and the two exact cubic endpoint
signs. A generic interval-evaluation kernel lifts those lower bounds
to the actual root. A finite positive-basis kernel proves the
nonnegative-sum/strict-constant-coefficient inference.

\paragraph{Complete entrance and unbounded induction.}
The 188 pairs $1\le u+r\le18$, excluding ONLY (0,2), are a
complete small entrance. Their 3228 retained coefficients and
2856 adjacent minors are explicitly positive; the D=1 empty prefix
is included as a vacuous context. The independent checker rebuilds
EVERY polynomial by closed binomial/trinomial sums and computes its
inverse in the different basis $\mathbb Q(a)$, with no iterated
product or borrowed primary inverse. The same complete scalar
certificate verifies all these signs.

For every next w>=19 choose a predecessor of weight w-1 by removing
a ternary factor when r>0, and otherwise a binary factor. Its count
is at least eighteen, so it is not the single excluded pair.
Positive CI convolution covers all inherited ranks, and
(2)--(5) cover all new values and minors. Induction proves the theorem
for EVERY count. A complete natural-index predicate induction is
included among the five generic closure reports.

At (u,r)=(0,2), the retained prefix is
$(1,2-a,4-3a)$; its centre-one minor is
$a^2-a=(\lambda-1)(\lambda-2)<0$. The exact exception is therefore
necessary, not removed by silently strengthening the entrance.

\paragraph{Scope.}
The thirteen new scoped reports comprise six terminal-norm kernels,
five root/interval/positive-basis/all-count/exception kernels, and
two aggregate finite scalar/root-sign certificates. Ordinary LC
convolution, coefficient identifications, the independently certified
polynomial identities and the full sequence theorem are the written
assembly above. There is no end-to-end Lean composition datatype
formalization. The finite entrance and all polynomial certificates
are explicit mathematical premises verified in full.
This is a new finite-kernel theorem for the TRUE tribonacci root
factor; its transfer to full actual difference profiles has not yet
been proved. Earlier actual C, full A2, the u>=33 actual diagonal and
ORIGINAL undilated coefficient unimodality remain OPEN.
